\Lecturer{Hui Jin}
\Contributor{}
\LectureNumber{15}
\LectureDate{DATE}
\LectureTitle{CSS Codes and Steane Code}

\lstset{style=mystyle}


	\MakeScribeTop

\section{Stabilizer Codes Introduction}

A \emph{stabilizer code} is a quantum error-correcting code described by an abelian subgroup
\(\mathcal{S}\) of the (n)-qubit Pauli group \(\mathcal{P}_n\). The code space is the joint
(+1)-eigenspace of all operators in \(\mathcal{S}\).

\subsection*{Pauli group}
The (n)-qubit Pauli group \(\mathcal{P}_n\) is generated by tensor products of the single-qubit operators
\[
I = \begin{pmatrix}1 & 0 \ 0 & 1\end{pmatrix}, \quad
X = \begin{pmatrix}0 & 1 \ 1 & 0\end{pmatrix}, \quad
Y = \begin{pmatrix}0 & -i \ i & 0\end{pmatrix}, \quad
Z = \begin{pmatrix}1 & 0 \ 0 & -1\end{pmatrix},
\]
together with global phases \({\pm 1, \pm i}\).

\subsection*{Stabilizer group}
A stabilizer code is defined by an abelian subgroup \(\mathcal{S} \subset \mathcal{P}_n\), not containing \(-I\).
Each \(S \in \mathcal{S}\) imposes the constraint
\[
S , |\psi\rangle = |\psi\rangle, \quad \forall S \in \mathcal{S}.
\]

If \(\mathcal{S}\) is generated by \(n-k\) independent elements, then the code space has dimension \(2^k\). Such a code is called an \([[n,k,d]]\) code.

\subsection*{Error detection and correction}
Errors are modeled as Pauli operators \(E \in \mathcal{P}_n\).
\begin{itemize}
\item If \(E\) commutes with all stabilizers, it preserves the code space and acts as a logical operator.
\item If \(E\) anticommutes with some stabilizer, measurement of that stabilizer yields eigenvalue \(-1\).
\end{itemize}

The pattern of \(\pm 1\) outcomes is called the \emph{syndrome}, which identifies the error up to equivalence.

The minimum weight of a Pauli operator that maps one valid codeword to another (without being in \(\mathcal{S})\) is the \emph{distance} \(d\). An \([[n,k,d]]\) stabilizer code corrects up to \(\lfloor (d-1)/2 \rfloor\) errors.

\subsection*{CSS construction}
A special class of stabilizer codes, due to Calderbank--Shor and Steane (CSS codes), arises from pairs of classical binary linear codes \((C_X, C_Z)\) with
\[
C_Z^\perp \subset C_X.
\]

Here,
\begin{itemize}
\item \(Z\)-type stabilizers are defined from parity checks of \(C_Z\),
\item \(X\)-type stabilizers from parity checks of \(C_X\).
\end{itemize}

The Steane code \([[7,1,3]]\) is a canonical example of this construction.


\section{Steane code in Stabilizer Formulation}
In the CSS construction of the Steane code, each row \(h \in H\) of the parity-check matrix defines both a \(Z\)-type stabilizer \(Z^h\) and an \(X\)-type stabilizer \(X^h\).

\[
X^h = \bigotimes_{i=1}^7 X^{h_i},
\qquad
Z^h = \bigotimes_{i=1}^7 Z^{h_i},
\]
where \(h=(h_1,\dots,h_7)\).

---

\subsection{Action on basis states}

For a computational basis state \(|v\rangle\), with \(v\in {0,1}^7\),
\[
X^h |v\rangle = |v \oplus h\rangle,
\]
where \(\oplus\) is bitwise XOR.

Thus, an \(X\)-stabilizer does not leave \(|v\rangle\) invariant; it maps it to another basis state shifted by \(h\).

---

\subsection{Stabilizer condition}

If a state \(|\psi\rangle\) is stabilized by \(X^h\),
\[
X^h |\psi\rangle = |\psi\rangle,
\]
then whenever \(|v\rangle\) has nonzero amplitude in \(|\psi\rangle\), the shifted state \(|v\oplus h\rangle\) must also appear with the same amplitude. Hence, the amplitudes are constant on orbits generated by the shifts \(h\).


\subsection{Consequence for CSS codes}

\begin{itemize}
\item Z stabilizers: restrict the support to classical codewords \(C = \ker(H)\).
\item X stabilizers: enforce uniform superpositions across cosets generated by \(C^\perp\).
\end{itemize}


For the Steane code, where \(C_X = C_Z = C\) (the Hamming \([7,4,3]\) code), the result is:
\begin{itemize}
    \item \(|0_L\rangle\) is the equal superposition of the 8 codewords of \(C\),
    \item \(|1_L\rangle = X^{\otimes 7}|0_L\rangle\) is the equal superposition of the complementary coset.
\end{itemize}
Thus, the \(Z\)-part selects the codewords, while the \(X\)-part forces them to appear with equal weight.
